\chapter{Bermudans and Callables}\label{ch:rc:bermudans-and-callables}

In 2014 Taiwan's legislature exempted foreign-currency bonds listed on the
island from the cap on its life insurers' overseas investments, and the
insurers, with hundreds of billions of dollars to place, began buying thirty-year
dollar zero-coupon bonds that the issuing banks could call every year after
the first few. The investors earned a yield well above the Treasury curve for
selling that right. The issuers did not keep it: they swapped each bond into
floating funding with a dealer, and the swap carried the call, so that dealers
ended up long a large stock of Bermudan options on long-dated rates. To hedge,
they sold volatility in the swaption market, year after year. The instrument at
the centre is the \omterm{def:rc:bermudans-and-callables:bermudan}{Bermudan swaption}: an option to enter a swap on any one of
several dates. It is worth more than the most valuable of the Europeans it can
become, by an amount no European hedge captures. This chapter prices it on
chapter~7's tree and by regression Monte Carlo, measures that amount, and prices
the callable zero-coupon bond the insurers bought.

\section{Bermudan swaptions and their co-terminal Europeans}

\begin{definition}[Bermudan swaption]\label{def:rc:bermudans-and-callables:bermudan}
A \emph{Bermudan swaption}\index{Bermudan swaption} gives the right, on any one of
a set of exercise dates $T_{e_1}<\dots<T_{e_m}$, to enter the swap that starts on
that date and ends on a fixed final date $T_n$, at a fixed rate $K$. A
``ten-non-call-one'' receiver Bermudan can be exercised every year from year one
into a receiver swap ending at year ten. It is a Bermudan exercise right in the
sense of One Quant Book 5, chapter 6, applied to swaps.
\end{definition}

\begin{definition}[Co-terminal swaption]\label{def:rc:bermudans-and-callables:coterminal}
The \emph{co-terminal swaptions}\index{co-terminal swaption} of a Bermudan are the
European swaptions expiring on each of its exercise dates into the swap ending on
its final date: the Europeans it can become.
\end{definition}

\begin{definition}[Switch option]\label{def:rc:bermudans-and-callables:switch}
The \emph{switch option}\index{switch option} of a Bermudan is its value minus the
value of its most expensive co-terminal European: the value of being able to
choose the exercise date after seeing the rates.
\end{definition}

\begin{proposition}[Bounds from the Europeans]\label{prop:rc:bermudans-and-callables:bounds}
A Bermudan is worth at least its most expensive co-terminal European and at most
the sum of all of them.
\end{proposition}
\begin{proof}
The holder can commit to one exercise date, which gives the European; any
exercise policy exercises on at most one date, so the payoff is dominated by the
sum of all the Europeans' payoffs.
\end{proof}

The lower bound is what a static hedge with Europeans secures; the \omterm{def:rc:bermudans-and-callables:switch}{switch option}
depends on how the co-terminal swap rates move together, which only a model of
the whole curve can price.

\section{Exercise by backward induction on a tree}

\begin{method}[A Bermudan on the Hull--White tree]\label{met:rc:bermudans-and-callables:tree}
Build chapter~7's tree to the final date. Roll back from the last exercise date:
at each exercise date, compute at every node the value of the underlying swap
(its zero-coupon bonds rolled back on the tree), and replace the continuation
value by the maximum of the two; between exercise dates, discount along the
branches. The value at the root is the Bermudan's price; the nodes where
exercise wins trace the exercise boundary.
\end{method}

\omcode{code/firm/bermudan/firm_bermudan.py}{29}{41}{A Bermudan swaption by backward induction on the tree.}\label{lst:rc:bermudans-and-callables:tree}

\begin{example}[A ten-non-call-one receiver]\label{ex:rc:bermudans-and-callables:tenno}
On chapter~2's euro OIS curve, with Hull--White at $\kappa=3\%$ and $\sigma=86$ basis
points (chapter~7's single-volatility fit), a receiver Bermudan at the ten-year par
rate of 2.585\% exercisable every year from one to nine is worth 358.4 basis points
of notional. Its co-terminal Europeans are worth 209.5, 245.3, 249.5, 237.8,
214.1, 184.3, 147.8, 104.3 and 54.9 basis points (tree, within 1.4 of Jamshidian's
formula; \cref{fig:rc:bermudans-and-callables:europeans}); the most expensive is
the three-year. The \omterm{def:rc:bermudans-and-callables:switch}{switch option} is worth 109.0 basis points, 44\% of the best
European.
\end{example}

\begin{omfigure}
\begin{tikzpicture}
\begin{axis}[width=11.5cm,height=5.4cm,ybar,bar width=14pt,
  xlabel={exercise date (years); swap ends at year ten},ylabel={value (bp of notional)},
  tick label style={font=\small},label style={font=\small},
  xmin=0.4,xmax=9.6,ymin=0,ymax=420,xtick={1,2,3,4,5,6,7,8,9},grid=major,grid style={gray!25},
  table/col sep=comma]
\addplot[fill=omDef!55,draw=omDef] table[x=e,y=euro]{figdata/rates-credit-risk/09-bermudans-and-callables/europeans.csv};
\addplot[omThm,very thick,sharp plot,mark=none,forget plot] table[x=e,y=berm]{figdata/rates-credit-risk/09-bermudans-and-callables/europeans.csv};
\node[font=\footnotesize,text=omThm,above] at (axis cs:7,365) {Bermudan};
\end{axis}
\end{tikzpicture}

\omcaption{The co-terminal receiver swaptions of a ten-non-call-one euro Bermudan
at the ten-year par rate, and the Bermudan itself (the line). The gap between the
line and the tallest bar is the \omterm{def:rc:bermudans-and-callables:switch}{switch option}. Data: chapter~2's curve,
Hull--White with $\kappa=3\%$ and $\sigma=86$ basis points; the chapter's
tutorial.}\label{fig:rc:bermudans-and-callables:europeans}
\end{omfigure}

The exercise boundary of the receiver is a short rate below which exercising
beats waiting (\cref{fig:rc:bermudans-and-callables:boundary}). Early on the rate
must fall far, to 0.82\% in year one, since exercising gives up eight more
chances; by year nine exercise is worth it below 2.30\%, close to the strike.

\begin{omfigure}
\begin{tikzpicture}
\begin{axis}[width=11.5cm,height=5.2cm,
  xlabel={exercise date (years)},ylabel={rate (\%)},
  tick label style={font=\small},label style={font=\small},
  xmin=0.5,xmax=9.5,ymin=0.5,ymax=3.3,xtick={1,2,3,4,5,6,7,8,9},grid=major,grid style={gray!25},
  legend columns=3,legend style={font=\footnotesize,at={(0.5,-0.25)},anchor=north,draw=none,fill=none,/tikz/every even column/.append style={column sep=8pt}},
  table/col sep=comma]
\addplot[omThm,very thick,mark=*] table[x=e,y=rstar]{figdata/rates-credit-risk/09-bermudans-and-callables/boundary.csv};
\addplot[omDef,thick,dashed] table[x=e,y=strike]{figdata/rates-credit-risk/09-bermudans-and-callables/boundary.csv};
\addplot[omProp,thick,dotted] table[x=e,y=fwd]{figdata/rates-credit-risk/09-bermudans-and-callables/boundary.csv};
\legend{exercise below this short rate,strike 2.585\%,one-year forward rate}
\end{axis}
\end{tikzpicture}

\omcaption{Exercise boundary of the receiver Bermudan on the tree: the highest
short rate at which exercising into the remaining swap beats continuing. It rises
towards the strike as the exercise opportunities run out. Data: the chapter's
tutorial.}\label{fig:rc:bermudans-and-callables:boundary}
\end{omfigure}

\begin{example}[Mean reversion and the switch]\label{ex:rc:bermudans-and-callables:kappa}
Recalibrating $\sigma$ so that the five-into-five European keeps its price, the
Bermudan and its \omterm{def:rc:bermudans-and-callables:switch}{switch option} depend on $\kappa$ (table below): stronger mean
reversion decorrelates the rates of successive exercise dates and makes the choice
of date more valuable, while the best European barely moves.
\end{example}

\begin{center}
\small
\begin{tabular}{rrrrr}
\toprule
$\kappa$ & $\sigma$ (bp) & Bermudan (bp) & best European (bp) & switch (bp) \\
\midrule
1\% & 78.1 & 351.0 & 249.2 & 101.9 \\
3\% & 86.0 & 358.4 & 249.5 & 109.0 \\
6\% & 98.9 & 369.7 & 249.7 & 120.0 \\
\bottomrule
\end{tabular}

\footnotesize Mean reversion, recalibrated volatility and the ten-non-call-one
receiver. Data: the chapter's tutorial.
\end{center}

\section{Exercise by regression in Monte Carlo}

Trees serve one or two factors; a market model (chapter~8) or any model with
many state variables is simulated, and simulation runs forward while exercise is
decided backward. The Longstaff--Schwartz method (One Quant Book 5, chapter 23)
estimates the continuation value by regressing, at each exercise date, the
discounted future cash flows of in-the-money paths on functions of the state.

\begin{method}[Regression exercise for a Bermudan swaption]\label{met:rc:bermudans-and-callables:lsm}
Simulate the state under a convenient measure (here the terminal forward measure
of year ten, numeraire $P(t,T_n)$) at the exercise dates. From the last date back:
compute the exercise value in numeraire units; on in-the-money paths regress the
realised future cash flow on $1$, the exercise value and its square; exercise
where the exercise value beats the fitted continuation. Fit the regressions on one
set of paths and apply the rule to an independent set: the result is a price of a
feasible policy, hence a lower bound.
\end{method}

\omcode{code/firm/bermudan/firm_bermudan.py}{121}{142}{The regression step of Longstaff--Schwartz, run backwards through the exercise dates.}\label{lst:rc:bermudans-and-callables:lsm}

\begin{example}[Tree and regression agree]\label{ex:rc:bermudans-and-callables:lsm}
With 40\,000 paths under the ten-year forward measure (antithetic, half-monthly
steps) the regression price of the ten-non-call-one is 358.1 basis points with a
standard error of 1.9, against 358.4 on the tree.
\end{example}

\section{Upper bounds and the switch option}

A regression price is a lower bound; how far below the true value it lies is
measured by a dual upper bound (One Quant Book 5, chapter 23), which subtracts a
martingale from the payoff and takes the maximum over dates along each path:
the gap between the two bounds certifies the policy, and computing it (by nested
simulation) belongs in the pricer's validation. Once the gap is small, the
practical risks lie elsewhere: in the model's correlation and volatility
structure, which set the \omterm{def:rc:bermudans-and-callables:switch}{switch option}'s size.

\begin{remark}[Hedging a Bermudan]\label{rem:rc:bermudans-and-callables:hedge}
A Bermudan's vega sits on its \omterm{def:rc:bermudans-and-callables:coterminal}{co-terminal swaptions}, split roughly in proportion
to the probability of exercising into each; its \omterm{def:rc:bermudans-and-callables:switch}{switch option} is a bet on the
decorrelation of successive swap rates that no European captures. Desks hedge the
Europeans' vega and carry, with a reserve, the rest; a change of model (chapter~7's
$\sigma(t)$, a market model's $\beta$) moves the switch value.
\end{remark}

\section{The callable-bond issuance business}

\begin{definition}[Formosa bond]\label{def:rc:bermudans-and-callables:formosa}
A \emph{Formosa bond}\index{Formosa bond} is a bond denominated in a foreign
currency and issued and listed in Taiwan. From 2014 most dollar Formosas were
long-dated (twenty or thirty years) and callable by the issuer, many with zero
coupons, sold to Taiwanese life insurers.
\end{definition}

\begin{definition}[Callable range accrual]\label{def:rc:bermudans-and-callables:rangeaccrual}
A \emph{callable range accrual}\index{callable range accrual} is a callable note
whose coupon accrues only on days when a reference rate stays within a range; the
issuer's call makes it a Bermudan on a portfolio of digital options, and it is
priced with the same exercise methods, on a model that also fits the smile.
\end{definition}

\begin{omfigure}
\begin{tikzpicture}[font=\small,
  box/.style={draw,thick,rounded corners=2pt,align=center,minimum height=1.1cm,minimum width=2.8cm,inner sep=4pt},
  arr/.style={-{Stealth[length=2.5mm]},thick}]
\node[box,draw=omDef,fill=omDef!8] (inv) at (0,0) {life insurer\\(investor)};
\node[box,draw=omThm,fill=omThm!8] (iss) at (4.8,0) {issuing bank};
\node[box,draw=omProp,fill=omProp!8] (dlr) at (9.6,0) {dealer's\\rates desk};
\node[box,draw=gray,fill=gray!8] (mkt) at (9.6,-2.6) {swaption market};
\draw[arr] (inv.15) -- node[above=4pt,align=center] {cash at issue} (iss.165);
\draw[arr] (iss.195) -- node[below=4pt,align=center] {zero coupon, callable} (inv.345);
\draw[arr] (iss.15) -- node[above=4pt,align=center] {swap + call right} (dlr.165);
\draw[arr] (dlr.195) -- node[below=4pt,align=center] {floating funding} (iss.345);
\draw[arr,omExo] (dlr) -- node[right=4pt,align=left] {sells volatility\\to hedge the Bermudan} (mkt);
\end{tikzpicture}

\omcaption{The flows of a callable Formosa issue. The investor sells the call
right to the issuer for a higher yield; the issuer passes it to a dealer inside
the swap that turns the bond into floating funding; the dealer, long a Bermudan,
sells volatility to hedge it.}\label{fig:rc:bermudans-and-callables:flows}
\end{omfigure}

\begin{dated}{2026-09}{The Formosa market}\label{dat:rc:bermudans-and-callables:formosa}
In May 2014 Taiwan's legislature excluded locally issued foreign-currency bonds
from the 45\% cap on insurers' overseas investments, unlocking funds reported at
USD~586 billion; the dollar notes sold in the following months were all callable,
with maturities of twenty or thirty years, and 74\% paid no coupon (Bloomberg
data). The trade press has described the resulting hedging flows as a steady
supply of long-dated dollar volatility that helped hold it down.
\end{dated}

\section{Tutorial: a Bermudan three ways}\label{tut:rc:bermudans-and-callables:tutorial}

\begin{tutorial}
\textbf{Goal.} Price a ten-non-call-one receiver on the tree and by regression,
measure its \omterm{def:rc:bermudans-and-callables:switch}{switch option} and its exercise boundary. \textbf{End state:}
\cref{fig:rc:bermudans-and-callables:europeans,fig:rc:bermudans-and-callables:boundary}
and the numbers of \cref{ex:rc:bermudans-and-callables:tenno,ex:rc:bermudans-and-callables:lsm}.
\begin{tutsteps}
\item \textbf{Tree}: \texttt{bermudan\_summary()} builds chapter~7's tree with
half-monthly steps and prices the Bermudan and its nine Europeans.
\item \textbf{Boundary}: \texttt{boundary()}.
\item \textbf{Regression}: \texttt{lsm\_check()} on 40\,000 paths.
\item \textbf{Mean reversion}: \texttt{mean\_reversion\_table()};
\texttt{fig\_rc\_bermudan.py} writes the charts.
\end{tutsteps}
\textbf{What to change next.} Add $x^2$ and the one-year rate to the regression
basis and see whether the price rises; price the Bermudan on chapter~7's
piecewise-$\sigma$ model with Monte Carlo.
\end{tutorial}

\section{Build: Bermudan and callable pricers}\label{bld:rc:bermudans-and-callables:bermudan}

\begin{build}
\bfield{Purpose} The book's callable products: \omterm{def:rc:bermudans-and-callables:bermudan}{Bermudan swaptions}, callable bonds
and notes, and the exercise decisions of chapter~12's mortgages.
\bfield{Interface} \texttt{bermudan\_tree(tree, exercises, final, K, receiver)};
\texttt{exercise\_boundary}; \texttt{callable\_zero\_tree(tree, maturity,
accretion, calls)}; \texttt{simulate\_x}, \texttt{bermudan\_lsm(model, exercises,
final, K, paths)}; \texttt{switch\_option}.
\bfield{Rules} Chapter~7's \texttt{firm\_shortrate} imported; regression fitted and
applied on independent path sets; fixed seeds; numpy only.
\bfield{Acceptance tests} \texttt{code/firm/bermudan/tests/}: a single exercise
date gives the European; the Bermudan lies between the best European and the
sum; regression agrees with the tree; a zero accreting far above rates is
called at the first date and one accreting at zero almost never; the
boundary lies below the strike.
\bfield{Stretch} A dual upper bound by nested simulation; Bermudans in the market
model of chapter~8; \omterm{def:rc:bermudans-and-callables:rangeaccrual}{callable range accruals}.
\end{build}

\begin{omsources}
\item F.\ Longstaff and E.\ Schwartz, ``Valuing American options by simulation: a
simple least-squares approach'', \emph{Review of Financial Studies} 14(1), 2001.
\item L.\ Andersen and M.\ Broadie, ``Primal-dual simulation algorithm for pricing
multidimensional American options'', \emph{Management Science} 50(9), 2004.
\item Insurance Journal (Bloomberg News), ``Goldman leads Wall Street tapping
Taiwan cash pile'', 10 September 2014; Risk.net, ``Formosas, the Fed, and the
billion-dollar Bermudan trade''.
\item L.\ Andersen and V.\ Piterbarg, \emph{Interest Rate Modeling}, volume 3,
Atlantic Financial Press, 2010, chapter 19.
\end{omsources}

\section{Exercises}

\begin{exercise}[$\star$]\label{exo:rc:bermudans-and-callables:1}
Using \cref{ex:rc:bermudans-and-callables:tenno}, give the lower and upper bounds
of \cref{prop:rc:bermudans-and-callables:bounds} and check the Bermudan lies
between them.
\end{exercise}

\begin{exercise}[$\star$]\label{exo:rc:bermudans-and-callables:2}
Why is the three-year European the most expensive, and not the one-year with the
longest swap?
\end{exercise}

\begin{exercise}[$\star$]\label{exo:rc:bermudans-and-callables:3}
An investor buys a thirty-year callable zero instead of a straight one. What has
she sold, and to whom does it end up?
\end{exercise}

\begin{exercise}[$\star\star$]\label{exo:rc:bermudans-and-callables:4}
Read \cref{fig:rc:bermudans-and-callables:boundary}. Why does the boundary rise
towards the strike, and why does it stay below the one-year forward rate?
\end{exercise}

\begin{exercise}[$\star\star$]\label{exo:rc:bermudans-and-callables:5}
Explain, from the table of \cref{ex:rc:bermudans-and-callables:kappa}, why the \omterm{def:rc:bermudans-and-callables:switch}{switch
option} rises with $\kappa$ while the best European barely changes.
\end{exercise}

\begin{exercise}[$\star\star$]\label{exo:rc:bermudans-and-callables:6}
Why must the regression be fitted on one set of paths and applied to another to
give a lower bound?
\end{exercise}

\begin{exercise}[$\star\star\star$]\label{exo:rc:bermudans-and-callables:7}
\emph{Coding.} Price the thirty-non-call-five callable zero of the weekend problem
at its par accretion yield, and find the yield of the straight zero on the same
curve.
\end{exercise}

\begin{exercise}[$\star\star\star$]\label{exo:rc:bermudans-and-callables:8}
\emph{Find the flaw.} ``We hedge our Bermudan with the most expensive co-terminal
European, in the same notional, so we are hedged.''
\end{exercise}

\section{Problem: The Formosa Hedge}

\begin{problem}[{Weekend problem --- a thirty-year callable zero}]\label{pb:rc:bermudans-and-callables:1}
A bank issues a thirty-year dollar zero-coupon \omterm{def:rc:bermudans-and-callables:formosa}{Formosa bond}, callable every year from
year five at its accreted value, and sells it at par to a life insurer. It swaps the
bond with a dealer, who takes over the call. Price it with Hull--White ($\kappa=3\%$,
$\sigma=90$ basis points) on chapter~1's SOFR curve (monotone convex), monthly tree.

\medskip\noindent\textbf{Part I --- The bond.}
\begin{enumerate}
\item Give the thirty-year zero-coupon yield of the curve (annual compounding).
\item Give the accretion yield at which the callable zero is worth par.
\item Give the value of the straight zero at that accretion, per unit issued.
\item Give the value of the issuer's call.
\item Why does the investor accept this?
\end{enumerate}

\medskip\noindent\textbf{Part II --- The option.}
\begin{enumerate}[resume]
\item Which single call date is worth the most as a European, and how much?
\item Give the Bermudan premium over that European.
\item Why is the first call date the most valuable European here?
\item What is the dealer long, and against what?
\item Why does the dealer not simply keep the option?
\end{enumerate}

\medskip\noindent\textbf{Part III --- The hedge.}
\begin{enumerate}[resume]
\item Give the vega of the call per basis point of volatility, per USD~1 billion.
\item Which swaptions does the dealer sell to hedge it?
\item If ten such issues of USD~1 billion come in a month, what vega enters the market?
\item What happens to long-dated swaption volatility, and to the dealers' marks?
\item When rates fall and the bonds are called, what happens to the hedges?
\end{enumerate}

\medskip\noindent\textbf{Part IV --- Judgement.}
\begin{enumerate}[resume]
\item What risks does the insurer keep?
\item Why is the trade concentrated in few dealers, and why does that matter?
\item What model risk does the dealer carry?
\item State the \emph{named result}: the par accretion yield, the Bermudan premium over
the best European and the vega per USD~1 billion.
\item In one sentence: who sells volatility in this chain, and why?
\end{enumerate}
\end{problem}

\section{Interview questions}

\begin{interviewq}[$\star$ \iqroles{trader, bank}]\label{iq:rc:bermudans-and-callables:1}
What is a \omterm{def:rc:bermudans-and-callables:bermudan}{Bermudan swaption}, and why is it worth more than any of its Europeans?
\end{interviewq}

\begin{interviewq}[$\star\star$ \iqroles{researcher, developer}]\label{iq:rc:bermudans-and-callables:2}
Explain Longstaff--Schwartz for a \omterm{def:rc:bermudans-and-callables:bermudan}{Bermudan swaption}. Why is it a lower bound?
\end{interviewq}

\begin{interviewq}[$\star\star$ \iqroles{researcher}]\label{iq:rc:bermudans-and-callables:3}
Which parameter of a one-factor model drives the \omterm{def:rc:bermudans-and-callables:switch}{switch option}, and how?
\end{interviewq}

\begin{interviewq}[$\star\star$ \iqroles{trader}]\label{iq:rc:bermudans-and-callables:4}
Why do callable-bond issuance programmes affect long-dated swaption volatility?
\end{interviewq}

\begin{interviewq}[$\star\star\star$ \iqroles{researcher, risk}]\label{iq:rc:bermudans-and-callables:5}
How would you validate a Bermudan pricer?
\end{interviewq}

\begin{interviewq}[$\star\star\star$ \iqroles{developer}]\label{iq:rc:bermudans-and-callables:6}
Your regression Bermudan price moves by two basis points when you change the
seed. Is that acceptable, and what would you do?
\end{interviewq}
